\documentclass[11pt]{article}
\usepackage[a4paper,margin=24mm]{geometry}
\usepackage{amsmath,amssymb,amsthm,booktabs}
\newtheorem{lemma}{Lemma}\newtheorem{theorem}{Theorem}
\title{Doubling, truncation and fusion joins\\for the $\operatorname{PSL}_3(4)$ catalogue}
\author{L100 finite source-join component}\date{}
\begin{document}\maketitle
We use the fixed Halle preprint of Gottschalk--Leemans \cite{GL}; the
later journal edition remains uncollated. Its labels $2.1,\ldots,5.4$
are joined to the accepted 29 literal coset representatives by the
previous source-join component. The complete subgroup inputs, broad
primitive enumeration, full-automorphism fusion and palette calculations
retain their exact previously stated premises. In particular local
two-transitivity was tested, not imposed as an enumeration filter.

\begin{theorem}
The four printed doubling links are valid:
\[
2.7\longrightarrow2.6,\quad4.1\longrightarrow4.8,\quad
5.2\longrightarrow5.3,\quad5.1\longrightarrow5.4.
\]
Of the 24 printed truncation links, 23 have complete literal incidence
maps. The remaining assertion, on physical page 9, that $3.6$ is a
truncation of $5.3$ is false by type cardinalities; replacing $5.3$ by
$5.2$ gives a verified truncation. All 29 printed parenthetical fusion
multiplicities agree with the complete acquired 145-family partition.
\end{theorem}
The truncation error concerns an ancillary link. It does not refute
the principal classification. This component does not identify the
abstract names of parabolics, kernels, automorphism or correlation groups.

\section{A concrete doubling argument}
\begin{lemma}
Let $\Gamma$ be a flag-transitive geometry with point type $X$ and
another type $E$. Suppose the $X$-shadows of the objects of $E$ are
distinct two-element sets. Assume that, for every object $z$ of any
other type,
\[
       e*z\quad\Longleftrightarrow\quad \sigma_X(e)\subseteq\sigma_X(z).
       \tag{1}
\]
Replace $X,E$ by copies $X^+,X^-$ of $X$. Declare $x^+*y^-$ when
$\{x,y\}=\sigma_X(e)$ for some $e\in E$; declare either copy of $x$
incident with $z$ exactly when $x*z$ in $\Gamma$; leave all other
incidences unchanged. The resulting incidence system is a geometry,
and the induced specified group action is flag-transitive.
\end{lemma}
\begin{proof}
A flag using both copies gives distinct points $x,y$ and a unique
object $e$ with shadow $\{x,y\}$. By (1), replacing $y^-$ by $e$
gives an old flag through $x,e$ and the unchanged objects. Extend it
to an old chamber. Its associated new chamber extends the original
flag. A flag using only $X^+$ extends from the corresponding old flag
through $x$; the other point of its chamber's $E$-object supplies the
minus copy. With only $X^-$ interchange these roles. With neither,
extend the unchanged old flag first. Thus every new flag extends.

The chamber bijection sends $(x,e,Z)$ to $(x^+,y^-,Z)$.
Here $y$ is the other point of $e$, and the bijection is equivariant.
Consequently old chamber transitivity implies new chamber transitivity;
in a geometry this implies flag transitivity by extending flags to
chambers. The proof works in rank two as well.
\end{proof}
Condition (1) also implies the particular $(I_{12})$ condition printed
on page 6. For an object $e$ and an object $z$ of another type, their
point-shadow intersection is empty, a singleton $\{x\}$, or the full
shadow of $e$. In the singleton case the flag $\{x\}$ is incident
with both and has the required point shadow. In the full case (1)
gives $e*z$, so the flag $\{e\}$ works. Incidence here is reflexive,
as in the source's coset definition.

The source's quoted Corollary 2.1 is explicitly a rank-at-least-three
statement attributed to its reference [40]. Its use at $2.7\to2.6$
therefore requires a qualification. The direct rank-two argument above
supplies that case: $2.7$ is the edge-incidence system of a graph, and
the construction is its bipartite double. We do not extend the scope of
the quoted corollary by assertion.

For $4.1,5.1,5.2$, the accepted projective model supplies ordinary
points and unordered point pairs. Incidence with a hyperoval means
containment, and incidence with the relevant line means disjointness,
for both the point and pair types. Thus (1) also has this immediate
projective explanation. Independently, the finite verifier checks every
pair shadow, its uniqueness, and every instance of (1) in the actual
coset models. No source point or line label is assumed to equal a
coset label.

\section{Exact identifications and all truncations}
All new maps use the same positive certificate interface. If $H_i$ is
an actual target parabolic, the certificate gives a derived source
object $x_i$ fixed by its images under an explicitly checked group
automorphism $\alpha$. On target right cosets the map is
\[
                         H_i g\longmapsto x_i^{\alpha(g)}.
\]
The derived geometry is either a doubled model from the lemma or an
induced system on a specified subset of source types. The original
coset actions are checked against the saved generator permutations;
right multiplication, not a guessed handedness, defines all object
images. Allowed automorphisms are literal conjugations by acquired
point normalizers and the acquired coordinate duality. Their exact
generator images and full group bijections are retained.

For every map the verifier checks all coset members, all typewise
bijections, every incidence and nonincidence row, and every generator
action on every object. Thus the target's full geometry and residue
properties transport, including its actual local two-transitivity
flags. This argument does not infer a map from matching group orders,
diagrams or object counts. The four doubling maps require respectively
$40320,80640,100800,100800$ coset-member checks.

The complete positive truncation ledger is as follows. Eight links
were already accepted in the projective-construction component and are
acquired unchanged; sixteen additional maps are checked here.
\begin{center}\small
\begin{tabular}{ll}\toprule
Target & Verified parent geometries\\\midrule
2.2 &5.1, 5.2, 4.10, 3.1, 3.4, 3.6\\
2.4 &4.11, 3.5\\
2.5 &3.3\\
3.1 &5.1, 5.2, 4.1\\
3.2 &4.3, 4.4\\
3.6 &5.1, \textbf{5.2}, 4.1, 4.10\\
4.1 &5.1, 5.2\\
4.8 &5.3, 5.4\\
4.9 &5.2\\
4.10&5.1\\\bottomrule
\end{tabular}
\end{center}
The bold entry repairs the printed $5.3$. Indeed, a truncation retains
the complete object sets of selected types. The type sizes of $3.6$
are $56,210,56$, whereas those of $5.3$ are $21,21,21,56,56$; these
are also visible directly in the native diagrams on pages 9 and 13.
No truncation of $5.3$ can introduce a 210-object type. The replacement
$5.2$ is established positively by its saved map. It is also consistent
with the projective constructions: $5.1$ and $5.2$ differ only in a
line--hyperoval relation, whereas this truncation keeps the two
hyperoval types and the point-pair type.

\section{Fusion multiplicities and evidence boundary}
The acquired complete fusion records contain 145 distinct inner-conjugacy
families with type permutations allowed, partitioned by the actual full
$\operatorname{Aut}(G)$ action. The rank totals are $21,25,63,36$ for
ranks $2,3,4,5$. Each retained family occurs once. The 29 representative
identities agree exactly with the independently acquired bare catalogue.
Grouping the saved membership records by these exact identities gives
the following source-order multiplicities:
\begin{center}\small
\begin{tabular}{cl}\toprule
Rank&Multiplicities in increasing source-label order\\\midrule
2&1,6,3,3,6,1,1\\
3&6,3,3,3,3,6,1\\
4&6,3,3,6,6,6,3,6,6,12,6\\
5&12,6,6,12\\\bottomrule
\end{tabular}
\end{center}
These agree with all parenthetical numbers on pages 7--13. The families
are counted modulo $G$-conjugacy and type permutations, then fused under
$\operatorname{Aut}(G)$, not extra bare geometries or correlation
automorphisms. The projection parser accepts only the complete
restricted membership-record syntax; it checks all 145 distinct
member identities and all 29 representative joins. The completeness
and inequivalence of those records are acquired from the accepted
enumeration and fusion proofs, not re-proved by counting their rows.

The first truncation pass saved ten complete positive maps before
rejecting the impossible printed link. That failed terminal receipt
is preserved. A changed pass handled only the six outstanding positive
maps, including the explicit replacement. Default replay verifies all
twenty new maps directly and does not repeat candidate-map discovery,
an ambient census or a graph-isomorphism search. Small-model controls
check the doubling construction and chamber correspondence; separate
damaged maps and record controls enforce the exact interfaces.
Abstract structural names, ordered-butterfly and oval constructions,
and the final whole-source synthesis remain separate. There is no new
measurement, Lean, question or priority claim.

\begin{thebibliography}{1}
\bibitem{GL} H. Gottschalk and D. Leemans,
\emph{Geometries for the group $\operatorname{PSL}(3,4)$},
Halle Reports on Algebra and Geometry 99--30 (1999), fixed 22-page preprint.
\end{thebibliography}
\end{document}
